\documentclass[11pt,a4paper]{article}
\usepackage[T1]{fontenc}
\usepackage[utf8]{inputenc}
\usepackage{amsmath,amssymb,mathtools}
\usepackage{booktabs}
\usepackage[hidelinks]{hyperref}
\usepackage[a4paper,margin=2.7cm]{geometry}

\newcommand{\R}{\mathbb{R}}
\newcommand{\trans}{\mathsf{T}}
\newcommand{\norm}[1]{\left\lVert #1 \right\rVert}
\newcommand{\code}[1]{\texttt{\detokenize{#1}}}

\title{Matrix Multiplication with Strassen's Algorithm\\
\large Theory and implementation in the \code{Strassen} example}
\author{PACS examples}
\date{\today}

\begin{document}
\maketitle

\begin{abstract}
This note describes Strassen's algorithm for dense matrix multiplication and
relates it to the implementation in
\code{Examples/src/LinearAlgebra/Strassen}. The code is not intended to replace
Eigen's optimized product indiscriminately. Instead, it uses a hybrid
algorithm: Strassen recursion is applied only above a configurable threshold,
while Eigen computes products at the lower recursion levels.
\end{abstract}

\section{The matrix multiplication problem}

Given $A\in\R^{m\times k}$ and $B\in\R^{k\times n}$, we seek to compute
$C=AB\in\R^{m\times n}$. For square $n\times n$ matrices, the classical
algorithm performs approximately $n^3$ scalar multiplications and a comparable
number of additions. Its arithmetic complexity is therefore $O(n^3)$.

Strassen observed that splitting each square matrix into four equally sized
blocks permits seven block products instead of the eight required by ordinary
block multiplication. Recursing on these products gives
\[
 T(n)=7T(n/2)+O(n^2),
\]
and hence
\[
 T(n)=O\left(n^{\log_2 7}\right)\approx O(n^{2.807}).
\]
This is an asymptotic improvement; it does not necessarily yield a lower
execution time for finite matrix sizes. Strassen requires additional additions,
temporary storage, and less regular memory-access patterns.

\section{Strassen's formula}

For simplicity, suppose that $n$ is even and write
\[
 A=\begin{bmatrix}A_{11}&A_{12}\\A_{21}&A_{22}\end{bmatrix},
 \qquad
 B=\begin{bmatrix}B_{11}&B_{12}\\B_{21}&B_{22}\end{bmatrix},
\]
where each block has size $(n/2)\times(n/2)$. Instead of calculating the
eight products of ordinary block multiplication, Strassen forms
\begin{align*}
 M_1 &= (A_{11}+A_{22})(B_{11}+B_{22}), &
 M_2 &= (A_{21}+A_{22})B_{11}, \\
 M_3 &= A_{11}(B_{12}-B_{22}), &
 M_4 &= A_{22}(B_{21}-B_{11}), \\
 M_5 &= (A_{11}+A_{12})B_{22}, &
 M_6 &= (A_{21}-A_{11})(B_{11}+B_{12}), \\
 M_7 &= (A_{12}-A_{22})(B_{21}+B_{22}).
\end{align*}
The result blocks are reconstructed as
\begin{align*}
 C_{11} &= M_1+M_4-M_5+M_7, &
 C_{12} &= M_3+M_5, \\
 C_{21} &= M_2+M_4, &
 C_{22} &= M_1-M_2+M_3+M_6.
\end{align*}
Only the seven products $M_i$ are recursive operations; all other operations
are block additions or subtractions, with local $O(n^2)$ cost.

\section{Why the algorithm is hybrid}

Recursing down to very small matrices is normally inefficient. The constant
costs of additions, recursive calls, temporaries, and weaker memory locality
can dominate the savings in multiplications. Moreover, Eigen implements
blocked and vectorized dense products that are highly efficient at moderate
sizes.

For this reason, the public interface
\[
  \code{apsc::strassen(A, B, cutoff)}
\]
uses Strassen only when the problem size exceeds \code{cutoff}; below the
threshold it directly executes Eigen's product
\[
  C=AB.
\]
The best threshold depends on the architecture, compiler, scalar type, problem
size, and Eigen or BLAS implementation. There is no universally optimal value:
the benchmark included in this directory is designed to measure it.

\section{Mapping to the code}

\subsection{Interface and dispatcher}

The template function \code{apsc::strassen} in \code{strassen.hpp} accepts
compatible Eigen expressions, verifies $A.cols()=B.rows()$, and evaluates each
expression once into an owned dense matrix. This avoids both recomputation of
Eigen expressions and a proliferation of nested template types in the
recursion.

The \code{detail::strassen_dispatch} dispatcher handles:
\begin{itemize}
\item empty matrices, returning a zero matrix of the correct size;
\item problems below the threshold, using Eigen's direct GEMM;
\item square matrices of equal order, starting the recursion immediately;
\item rectangular products or square matrices of unequal order, padding once
      to a square matrix and cropping the useful result block afterwards.
\end{itemize}

The initial padding uses the even order
\[
  N=2\left\lceil
  \frac{\max\{m,k,n\}}{2}\right\rceil,
\]
filling the added portions with zeros. It makes square recursion possible, but
can increase work and memory usage, especially for highly rectangular products.

\subsection{Recursive kernel and product assembly}

\code{detail::strassen_impl} receives non-owning views of $A$, $B$, and the
destination $C$. For an even order, it creates the four quadrant views, forms
the linear combinations required by $M_1,\ldots,M_7$, and recursively invokes
itself. The seven products are not all retained in memory: each product is
written to the current level's \code{M} buffer and immediately added to or
subtracted from the appropriate blocks of $C$. The accumulation order in the
code follows exactly the formula in the preceding section.

The base case is \code{gemm_into}, which executes
\code{dst.noalias() = A * B}. The \code{noalias()} declaration tells Eigen that
destination and operands do not overlap, allowing the kernel to avoid
unnecessary temporaries.

\subsection{Odd dimensions}

Strassen recursion requires division into equally sized quadrants. When a
square subproblem has odd order, this implementation does not pad it to the
next even order at every level. Instead, \code{strassen_peel} separates the
last row and column, recurses on the $(n-1)\times(n-1)$ block, and completes
the result using a rank-one update, two matrix-vector products, and one dot
product. This correction costs $O(n^2)$ and avoids repeated padding.

\subsection{Temporary-memory management}

The \code{ScratchPool} structure preallocates three temporary matrices per
level: \code{TA} and \code{TB} for operand linear combinations, and \code{M}
for the recursive product. Recursion is sequential, so sibling calls at the
same level reuse the same buffers. The pool is built once by the dispatcher and
the recursive kernel performs no heap allocations. Memory for square blocks is
bounded by a geometric series of order $4n^2/3$, in addition to small buffers
used by the odd-size treatment.

Eigen views such as \code{Ref}, \code{Block}, and \code{topLeftCorner} make it
possible to operate on quadrants without copying. This combination of views
and preallocated buffers makes the code closer to a practical numerical kernel
than a literal translation that allocates new temporaries at each recursive
call.

\section{Benchmarking and numerical correctness}

\code{main_StrassenBenchmark.cpp} compares \code{apsc::strassen} with Eigen's
direct product for a set of sizes and cutoffs. Each measurement is repeated,
sorted, and reported as minimum, median, and maximum elapsed time. The program
emits CSV fields for size, method, cutoff, timing, speedup, and relative error:
\[
  \frac{\norm{C_{\mathrm{Eigen}}-C_{\mathrm{Strassen}}}}
       {\norm{C_{\mathrm{Eigen}}}}.
\]
The reported speedup is the ratio of Eigen's median time to Strassen's median
time; a value larger than one favors Strassen.

The script \code{plot_strassen_benchmark.py} reads the CSV and creates Eigen
timing plots and heat maps for Strassen time, speedup, and relative error.
Because Strassen changes the order of additions and subtractions, floating
point results are not expected to be bitwise identical to Eigen's. The relative
error reveals the effect of these rounding differences.

\section{Limitations and extensions}

This example is intended for teaching and benchmarking. For general
applications, Eigen or an optimized BLAS library will often remain the better
choice. Possible extensions include automatic cutoff selection, parallel
kernels, a rectangular-matrix-specific memory strategy, and Strassen or
Winograd variants using fewer additions. Any performance comparison should be
run on the target system with data types and matrix sizes representative of the
application.

\begin{thebibliography}{9}
\bibitem{Strassen}
V.~Strassen, Gaussian elimination is not optimal,
\emph{Numerische Mathematik}, 13, pp.~354--356, 1969.

\bibitem{GolubVanLoan}
G.~H.~Golub and C.~F.~Van Loan,
\emph{Matrix Computations}, 4th ed., Johns Hopkins University Press, 2013.

\bibitem{Higham}
N.~J.~Higham, \emph{Accuracy and Stability of Numerical Algorithms},
2nd ed., SIAM, 2002.

\bibitem{IronyToledoTiskin}
D.~Irony, S.~Toledo, and A.~Tiskin,
Communication lower bounds for distributed-memory matrix multiplication,
\emph{Journal of Parallel and Distributed Computing}, 64(9), pp.~1017--1026,
2004.
\end{thebibliography}

\end{document}
